\chapter*{Abstract}
\addcontentsline{toc}{chapter}{Abstract}
\makeatletter\@mkboth{}{\scshape Abstract}\makeatother

\selectlanguage{english}

This project investigated the use of dynamic time warping (DTW) as a detection system for Bryde's whale calls in False Bay. Three 
feature extraction methods were considered and compared: the short-time Fourier transform (STFT), mel-frequency cepstral coefficients 
(MFCCs), with and 
without first and second derivatives, and linear predictive coding coefficients (LPCCs). Each feature extraction method was tested using 
both 5 templates per call type and 10 templates per call type. 

During hyperparameter tuning, the per-call-type thresholds, feature frame length and vote fraction for each model were selected. 
Selection decisions were made based on the highest f1-scores for each tested setting.
After tuning, call detection was performed on a full 72-hour audio file. The performance of each model was analysed using precision, recall, 
f1-score and average precision (AP).

The STFT model outperformed the other two models with a maximum f1-score of 0.792 and a maximum AP of 0.860. 
Precision was consistently lower than recall across all models, indicating that the models were more lenient during evaluation than intended in hyperparameter 
tuning.

Despite limitations caused by data scarcity and labelling inconsistencies, the DTW detection system was able to efficiently detect calls. This was 
especially true for the STFT feature extraction method.